\chapter{Применение формул аппроксимации для баллистического проектирования}

% =========================================================
% Рабочий каркас главы 4.
% Это НЕ финальная редакция главы, а удобная структура для
% аккуратного переноса материалов из ВКР, презентации и
% прикладных результатов по ILTT_Toolbox.
%
% Принцип заполнения:
% - сначала переносится опорный текст, таблицы, формулы и иллюстрации;
% - затем материал перестраивается по логике кандидатской диссертации;
% - служебные комментарии и указания на источники удаляются только
%   на позднем этапе шлифовки.
% =========================================================

% Общая опора:
% ВКР: глава 5, с. 60--70
% Презентация: слайды 33--38

\section{Сценарии использования полученных формул аппроксимации}
% Опора:
% ВКР: глава 5, с. 60--70
% Презентация: слайд 33
%
% Здесь ввести прикладной ракурс главы:
% - результаты глав 2 и 3 рассматриваются как инструмент
%   баллистического проектирования;
% - формулы могут использоваться не изолированно, а в составе
%   последовательной процедуры получения начального приближения;
% - необходимо развести основные сценарии использования:
%   1) оценка параметров перелёта;
%   2) построение начального приближения;
%   3) последующая оптимизация в параметрах начального приближения.
%
% Желательно сразу зафиксировать, что речь идёт не о замене
% численной оптимизации, а об ускорении и упрощении её запуска.
%
% Черновой текст:

\section{Алгоритм выбора начального приближения}
% Опора:
% ВКР: §5.1, с. 60--62
% Презентация: слайд 34
%
% Здесь последовательно описать полный алгоритм:
% 1) задаётся дата старта;
% 2) по формулам главы 2 оценивается продолжительность перелёта;
% 3) по формулам главы 2 восстанавливаются параметры управления
%    для перелёта между плоскими компланарными орбитами;
% 4) по кеплеровым элементам строятся матрицы линейных преобразований
%    из главы 3;
% 5) эти преобразования применяются к параметрам управления;
% 6) формируется начальное приближение для дальнейшей оптимизации.
%
% Важно:
% - текстом явно связать шаги главы 4 с материалом глав 2 и 3;
% - при переносе удобно сопровождать этот раздел одной
%   иллюстрацией-блок-схемой или примером перелёта.
%
% Черновой текст:

\section{Пример построения начального приближения для перелёта к Марсу}
% Опора:
% ВКР: §5.1, с. 60--62
% Презентация: слайд 34
%
% Здесь оставить место под подробный пример:
% - конкретная дата старта;
% - оценка продолжительности перелёта;
% - выбор параметров многообразия;
% - применение линейных преобразований;
% - сравнение начального приближения и оптимальной траектории.
%
% Этот раздел удобно делать максимально наглядным:
% - одна схема шагов;
% - одна иллюстрация траекторий;
% - краткие пояснения, почему начальное приближение уже достаточно
%   близко к оптимальному решению.
%
% Черновой текст:

\section{Ошибки начального приближения для перелётов Земля--Марс и Земля--Венера}
% Опора:
% ВКР: §5.2, с. 63--66
% Презентация: слайд 35
%
% Здесь собрать материал по невязкам начального приближения:
% - какие целевые параметры конечной орбиты контролируются;
% - каковы ошибки для сценариев Земля--Марс и Земля--Венера;
% - за какой временной диапазон / период они оценивались;
% - почему ошибки оказываются различными для разных планет.
%
% Важно:
% - сохранить таблицу ошибок как один из центральных прикладных
%   результатов главы;
% - отдельно проговорить, что деформация сильнее проявляется там,
%   где геометрия орбит дальше от плоской компланарной постановки.
%
% Черновой текст:

\section{Компенсация деформации и оптимизация в параметрах начального приближения}
% Опора:
% ВКР: §5.2, §5.3, с. 63--68
% Презентация: слайд 36
%
% Здесь нужно показать:
% - как компенсируется деформация траекторий;
% - почему удобнее оптимизировать не напрямую по сопряжённым
%   переменным, а по параметрам начального приближения;
% - что число таких параметров равно шести:
%   \lambda, \kappa, j, \Omega, \xi, \chi;
% - как это улучшает управляемость и интерпретируемость процедуры.
%
% Желательно подчеркнуть:
% - сама параметризация начального приближения уже является
%   самостоятельным практическим результатом работы.
%
% Черновой текст:

\section{Сравнение с прямой оптимизацией в сопряжённых переменных}
% Опора:
% ВКР: §5.2--5.3, с. 63--68
% Презентация: слайды 35--36 как контекст
%
% Здесь оставить место для важного сопоставления:
% - чем оптимизация в параметрах начального приближения удобнее
%   прямой работы в пространстве начальных сопряжённых переменных;
% - в чём состоит выигрыш с точки зрения устойчивости, интерпретируемости
%   и близости к физическим параметрам задачи;
% - как это соотносится с общей целью упрощения баллистического
%   проектирования.
%
% Если в окончательной версии отдельного формального сравнения не будет,
% этот раздел можно затем объединить с предыдущим.
%
% Черновой текст:

\section{Программная реализация в ILTT\_Toolbox}
% Опора:
% ВКР: глава 5, с. 60--70
% Презентация: слайды 36--37
%
% Здесь обязательно зафиксировать:
% - что предложенный инструментарий реализован программно;
% - название программы: Interplanetary Low Thrust Trajectories Toolbox;
% - что программа используется для работы с параметрами начального
%   приближения и задачами межпланетного перелёта;
% - что программа зарегистрирована как РИД.
%
% Здесь удобно отдельно развести:
% - научный результат диссертации;
% - его программную реализацию.
%
% При наличии точных реквизитов РИД их можно добавить позднее
% при заполнении библиографии и служебных данных.
%
% Черновой текст:

\section{Демонстрационный пример и область практической применимости}
% Опора:
% ВКР: глава 5, с. 60--70
% Презентация: слайды 37--38
%
% Здесь оставить место под:
% - демонстрацию работы ILTT_Toolbox;
% - примеры практического применения;
% - ограничения метода;
% - условия, в которых полученные формулы и начальные приближения
%   дают наибольший эффект.
%
% Важно:
% - этот раздел должен завершать прикладную линию диссертации;
% - здесь уместно кратко проговорить, где заканчивается область
%   надёжной применимости разработанного инструментария.
%
% Черновой текст:

\section{Применимость для других функционалов и направлений развития метода}
% Опора:
% ВКР: §5.4, с. 69--70
% Презентация: слайды 33--38 как общий контекст
%
% Здесь собрать:
% - насколько полученное начальное приближение применимо
%   за пределами исходного квадратичного функционала;
% - какие элементы подхода являются наиболее универсальными;
% - какие направления расширения метода выглядят естественными.
%
% Этот раздел удобно сделать мостом к заключению:
% - показать не только текущий результат, но и потенциал дальнейшего
%   развития.
%
% Черновой текст:

\section*{Выводы к главе 4}
% Опора:
% ВКР: глава 5, с. 60--70
% Презентация: слайд 38
%
% Здесь потом кратко зафиксировать:
% 1. прикладной сценарий использования формул аппроксимации
%    для баллистического проектирования;
% 2. алгоритм выбора начального приближения;
% 3. результаты оценки ошибок для перелётов Земля--Марс
%    и Земля--Венера;
% 4. переход к оптимизации в параметрах начального приближения;
% 5. программную реализацию метода в ILTT_Toolbox.
%
% Черновой текст:
